\documentclass[11pt]{article}
\usepackage{fontspec}
\setmainfont{Times New Roman}
\setsansfont{Arial}
\setmonofont{Consolas}
\usepackage{amsmath,amssymb,mathtools}
\usepackage{geometry}
\usepackage{microtype}
\geometry{a4paper,margin=2.35cm}
\setlength{\parindent}{1.25em}
\setlength{\parskip}{0pt}
\makeatletter
\renewcommand\footnotesize{\@setfontsize\footnotesize{7.7}{8.7}\selectfont}
\makeatother
\emergencystretch=100em
\allowdisplaybreaks
\sloppy
\hbadness=10000
\vbadness=10000
\hfuzz=2pt
\vfuzz=10pt
\raggedbottom

\begin{document}

% Unit: Paper31 Section04 entry 2 v001. Direct Interslavic Latin translation from the German final-audited source slice, source lines 245--267 / segment P31-S0057. Footnote counter continues after Paper31 Section04 entry 1.
\setcounter{footnote}{24}

1. \emph{Matrične kolca homomorfne k $\mathfrak R$.} Nehaj $\alpha_1,\ldots,\alpha_s$ bude linearno nezavisna $\mathfrak Z$-modulna baza ideala $\mathfrak a$, i nehaj $\gamma$ bude libovoljny element iz $\mathfrak R$. Togda $\gamma\alpha_i$ takože prinadleži k $\mathfrak a$, i imajut město ravnosti s koeficientami iz $\mathfrak Z$:
\[
        \gamma\alpha_i=c_{i1}\alpha_1+\cdots+c_{is}\alpha_s,
        \qquad i=1,\ldots,s.
\]
Ili v matričnoj formě, ako vsakokrat $(\tau_1,\ldots,\tau_s)$ označaje jednorednu matricu sostavljenu iz tyh elementov:
\[
 (\gamma\alpha_1,\ldots,\gamma\alpha_s)
 = (\alpha_1,\ldots,\alpha_s)
   \begin{pmatrix}
      c_{11}&\cdots&c_{1s}\\
      \vdots&&\vdots\\
      c_{s1}&\cdots&c_{ss}
   \end{pmatrix}
 = (\alpha_1,\ldots,\alpha_s)C.
\]
Kogda $\gamma$ proběgaje vse elementy iz $\mathfrak R$, cělost matric $C$, prirěđenyh pomoću bazy $\alpha_i$, tvori homomorfno kolco $R_\mathfrak a$; pri tom $R_\mathfrak a$ jest izomorfno k kolcu klasov ostatkov $\mathfrak R/((0):\mathfrak a)$, gde $(0):\mathfrak a$ označaje idealny kvocient nulovogo ideala črez $\mathfrak a$. Bo razlici $\beta-\gamma$ prirěđena jest razlika $B-C$, i to samo važi za produkt, poněže
\[
(\beta\gamma\alpha_1,\ldots,\beta\gamma\alpha_s)
=(\gamma\alpha_1,\ldots,\gamma\alpha_s)B
=(\alpha_1,\ldots,\alpha_s)CB.
\]
Elementam $c$ iz $\mathfrak Z$ osoblivo prirěđene sut diagonalne matrice $cE$. Vsim i samo tym elementam $\gamma$, za ktore $\gamma\mathfrak a$ jest ravno nulovomu idealu, prirěđena jest nulova matrica.

\end{document}
